\documentclass[../../../include/open-logic-section]{subfiles}

\begin{document}

\olfileid[hi]{sth}{story}{blv}

\olsection{परिशिष्ट: फ़्रेगे का मूल नियम V}

\olref[rus]{sec} में हमने बताया था कि रसेल ने अपने विरोधाभास को फ़्रेगे
द्वारा \emph{ग्रुंडगेज़ेत्से} में प्रतिपादित तंत्र की समस्या के रूप में
सूत्रित किया था। फ़्रेगे के तंत्र में सहज समुच्चय-निर्माण का सीधा सूत्रीकरण
नहीं था। अतः इस परिशिष्ट में हम अत्यंत संक्षेप में बताएँगे कि फ़्रेगे के तंत्र
में \emph{क्या} था, और उसका सहज समुच्चय-निर्माण तथा रसेल के विरोधाभास से
क्या संबंध है।

फ़्रेगे का तंत्र \emph{द्वितीय-क्रम} का है और उसे \emph{किसी संकल्पना के
विस्तार} की धारणा सूत्रित करने के लिए बनाया गया था।\footnote{ठीक-ठीक कहें
तो फ़्रेगे एक अधिक सामान्य धारणा, किसी फलन का ``मान-परास'', औपचारिक बनाना
चाहते हैं। संकल्पनाओं के विस्तार उस अधिक सामान्य धारणा के विशेष मामले हैं।
विवरण के लिए \citet[pp.\ 8--9]{Heck2012} देखें।} फ़्रेगे से प्रेरित संकेतन
का उपयोग करते हुए हम \emph{संकल्पना~$F$ के विस्तार} के लिए
$\fregeext{x}{F(x)}$ लिखेंगे। यह एक ऐसी युक्ति है जो \emph{विधेय} ``$F$''
को लेकर उसे (प्रथम-क्रम) \emph{पद} ``$\fregeext{x}{F(x)}$'' में बदलती है।
इस युक्ति से फ़्रेगे ने सदस्यता की निम्न \emph{परिभाषा} दी:
\[
	a \in b =_\text{df} \exists G(b = \fregeext{x}{G(x)} \land Ga)
\]
मोटे तौर पर: $a \in b$ तभी और केवल तभी जब $a$ ऐसी संकल्पना के अंतर्गत आता
है जिसका विस्तार $b$ है। (ध्यान दें कि परिमाणक ``$\exists G$'' द्वितीय-क्रम
का है।) फ़्रेगे ने निम्न सिद्धांत भी माना, जिसे \emph{मूल नियम V} कहते हैं:
$$\fregeext{x}{F(x)} = \fregeext{x}{G(x)} \liff \forall x (Fx \liff Gx)$$
मोटे तौर पर: संकल्पनाओं के विस्तार समान तभी और केवल तभी हैं जब वे
समविस्तारी हों। (फिर से, ``$F$'' और ``$G$'' दोनों विधेय-स्थान में हैं।) अब
एक सरल सिद्धांत सदस्यता को गुणधर्म-संतुष्टि से जोड़ता है:

\begin{lem}[\emph{ग्रुंडगेज़ेत्से} में]\ollabel{lem:Fregeextensions}
$\forall F \forall a(a \in \fregeext{x}{F(x)} \liff Fa)$
\end{lem}

\begin{proof} 
$F$ और $a$ को नियत करें। अब $a \in \fregeext{x}{F(x)}$ तभी और केवल तभी
जब $\exists G(\fregeext{x}{F(x)} = \fregeext{x}{G(x)} \land Ga)$ (सदस्यता
की परिभाषा से), तभी और केवल तभी जब
$\exists G(\forall x(Fx \liff Gx) \land Ga)$ (मूल नियम V से), तभी और केवल
तभी जब $Fa$ (प्राथमिक द्वितीय-क्रम तर्कशास्त्र से)।
\end{proof}

इससे सहज समुच्चय-निर्माण लगभग तुरंत मिलता है:

\begin{lem}[\emph{ग्रुंडगेज़ेत्से} में]
$\forall F \exists s \forall a (a \in s \liff Fa)$
\end{lem}

\begin{proof}
$F$ को नियत करें; अब \olref{lem:Fregeextensions} से $\forall a (a \in
\fregeext{x}{F(x)} \liff Fa)$ मिलता है; अतः अस्तित्वपरक सामान्यीकरण से
$\exists s\forall a(a \in s \liff Fa)$। चूँकि $F$ मनमाना था, परिणाम
निष्पन्न होता है।
\end{proof}

$F$ को $\forall x(Fx \liff x \notin x)$ द्वारा प्रदत्त मानने पर रसेल का
विरोधाभास निष्पन्न होता है।

\end{document}
